\chapter{Fluency bank}\label{AUTO}
Two questions per session, in 6 to 8 minutes including the correction. The identifiers allow revisits; they are not a list to be done in full at every lesson. The availability dates are based on the planned teaching, and should be shifted if the topic has not yet been studied.
\Needspace{5\baselineskip}\subsection*{R001 -- Fractions}\label{R001}
Calculate $\frac23+\frac14$.
\ifprof\textbf{Solution.} $\dfrac{11}{12}$: common denominator 12.\par\textit{Prerequisites:} Previous learning\par\textit{Available from 2026-09-03; suggested revisits:} 2026-09-10, 2026-09-24, 2026-11-03.\fi
\Needspace{5\baselineskip}\subsection*{R002 -- Order of operations}\label{R002}
Calculate $5-2(3-7)$.
\ifprof\textbf{Solution.} $13$: brackets, product, difference.\par\textit{Prerequisites:} Previous learning\par\textit{Available from 2026-09-03; suggested revisits:} 2026-09-10, 2026-09-24, 2026-11-03.\fi
\Needspace{5\baselineskip}\subsection*{R003 -- Powers}\label{R003}
Simplify $\dfrac{10^5\times 10^{-3}}{10^4}$.
\ifprof\textbf{Solution.} $10^{-2}=0.01$.\par\textit{Prerequisites:} Previous learning\par\textit{Available from 2026-09-03; suggested revisits:} 2026-09-10, 2026-09-24, 2026-11-03.\fi
\Needspace{5\baselineskip}\subsection*{R004 -- Roots}\label{R004}
Simplify $\sqrt{50}$.
\ifprof\textbf{Solution.} $5\sqrt2$.\par\textit{Prerequisites:} Previous learning\par\textit{Available from 2026-09-03; suggested revisits:} 2026-09-10, 2026-09-24, 2026-11-03.\fi
\Needspace{5\baselineskip}\subsection*{R005 -- Proportion}\label{R005}
18 represents 30 \% of a total. What is the total?
\ifprof\textbf{Solution.} $\dfrac{18}{0.3}=60$.\par\textit{Prerequisites:} Previous learning\par\textit{Available from 2026-09-03; suggested revisits:} 2026-09-10, 2026-09-24, 2026-11-03.\fi
\Needspace{5\baselineskip}\subsection*{R006 -- Direct rate}\label{R006}
A price of 80 falls by 15 \%. Final value?
\ifprof\textbf{Solution.} $80\times0.85=68$.\par\textit{Prerequisites:} Previous learning\par\textit{Available from 2026-09-03; suggested revisits:} 2026-09-10, 2026-09-24, 2026-11-03.\fi
\Needspace{5\baselineskip}\subsection*{R007 -- Initial value}\label{R007}
After a 20 \% increase, the value is 72. Initial value?
\ifprof\textbf{Solution.} $\dfrac{72}{1.2}=60$.\par\textit{Prerequisites:} Previous learning\par\textit{Available from 2026-09-03; suggested revisits:} 2026-09-10, 2026-09-24, 2026-11-03.\fi
\Needspace{5\baselineskip}\subsection*{R008 -- Calculated rate}\label{R008}
A quantity goes from 40 to 50. Rate?
\ifprof\textbf{Solution.} $\dfrac{50-40}{40}=25\,\%$.\par\textit{Prerequisites:} Previous learning\par\textit{Available from 2026-09-03; suggested revisits:} 2026-09-10, 2026-09-24, 2026-11-03.\fi
\Needspace{5\baselineskip}\subsection*{R009 -- Successive rates}\label{R009}
A 20 \% increase then a 20 \% decrease: overall rate?
\ifprof\textbf{Solution.} $1.2\times0.8=0.96$: a 4 \% decrease.\par\textit{Prerequisites:} Previous learning\par\textit{Available from 2026-09-03; suggested revisits:} 2026-09-10, 2026-09-24, 2026-11-03.\fi
\Needspace{5\baselineskip}\subsection*{R010 -- Reciprocal rate}\label{R010}
Which rate compensates a 20 \% decrease?
\ifprof\textbf{Solution.} $\dfrac{1}{0.8}=1.25$: a 25 \% increase.\par\textit{Prerequisites:} Previous learning\par\textit{Available from 2026-09-03; suggested revisits:} 2026-09-10, 2026-09-24, 2026-11-03.\fi
\Needspace{5\baselineskip}\subsection*{R011 -- Zero product}\label{R011}
Solve $(2x-1)(x+3)=0$.
\ifprof\textbf{Solution.} $x=\dfrac{1}{2}$ or $x=-3$.\par\textit{Prerequisites:} Previous learning\par\textit{Available from 2026-09-03; suggested revisits:} 2026-09-10, 2026-09-24, 2026-11-03.\fi
\Needspace{5\baselineskip}\subsection*{R012 -- Expansion}\label{R012}
Expand $(2x-3)^2$.
\ifprof\textbf{Solution.} $4x^2-12x+9$.\par\textit{Prerequisites:} Previous learning\par\textit{Available from 2026-09-03; suggested revisits:} 2026-09-10, 2026-09-24, 2026-11-03.\fi
\Needspace{5\baselineskip}\subsection*{R013 -- Factorisation}\label{R013}
Factorise $3x^2-12x$.
\ifprof\textbf{Solution.} $3x(x-4)$.\par\textit{Prerequisites:} Previous learning\par\textit{Available from 2026-09-03; suggested revisits:} 2026-09-10, 2026-09-24, 2026-11-03.\fi
\Needspace{5\baselineskip}\subsection*{R014 -- Sign of an affine function}\label{R014}
Solve $-3x+6\geq0$.
\ifprof\textbf{Solution.} $x\leq2$: reverse the inequality when dividing by -3.\par\textit{Prerequisites:} Previous learning\par\textit{Available from 2026-09-03; suggested revisits:} 2026-09-10, 2026-09-24, 2026-11-03.\fi
\Needspace{5\baselineskip}\subsection*{R015 -- Slope}\label{R015}
Slope of the line through $(1;4)$ and $(3;0)$?
\ifprof\textbf{Solution.} $-2$.\par\textit{Prerequisites:} Previous learning\par\textit{Available from 2026-09-03; suggested revisits:} 2026-09-10, 2026-09-24, 2026-11-03.\fi
\Needspace{5\baselineskip}\subsection*{R016 -- Line}\label{R016}
Equation of the line of slope -2 through $(1;4)$?
\ifprof\textbf{Solution.} $y=-2x+6$.\par\textit{Prerequisites:} Previous learning\par\textit{Available from 2026-09-03; suggested revisits:} 2026-09-10, 2026-09-24, 2026-11-03.\fi
\Needspace{5\baselineskip}\subsection*{R017 -- Weighted mean}\label{R017}
The values 2 and 5 have counts 3 and 2. Mean?
\ifprof\textbf{Solution.} $\dfrac{6+10}{5}=3.2$.\par\textit{Prerequisites:} Previous learning\par\textit{Available from 2026-09-03; suggested revisits:} 2026-09-10, 2026-09-24, 2026-11-03.\fi
\Needspace{5\baselineskip}\subsection*{R018 -- Median}\label{R018}
Median of 1, 2, 2, 5, 8, 10?
\ifprof\textbf{Solution.} $\dfrac{2+5}{2}=3.5$ using the convention of the mean of the two central values.\par\textit{Prerequisites:} Previous learning\par\textit{Available from 2026-09-03; suggested revisits:} 2026-09-10, 2026-09-24, 2026-11-03.\fi
\Needspace{5\baselineskip}\subsection*{R019 -- Quartiles}\label{R019}
For the data set 1, 2, 3, 4, 5, 6, 7, 8, give $Q_1$ and $Q_3$ with the ceiling ranks $\left\lceil\dfrac n4\right\rceil$ and $\left\lceil\dfrac{3n}4\right\rceil$.
\ifprof\textbf{Solution.} $Q_1=2$, $Q_3=6$.\par\textit{Prerequisites:} Previous learning\par\textit{Available from 2026-09-03; suggested revisits:} 2026-09-10, 2026-09-24, 2026-11-03.\fi
\Needspace{5\baselineskip}\subsection*{R020 -- Standard deviation of data}\label{R020}
For 1 and 3 with equal weights, mean, variance and standard deviation?
\ifprof\textbf{Solution.} Mean 2, variance $\dfrac{[(1-2)^2+(3-2)^2]}{2}=1$, standard deviation 1.\par\textit{Prerequisites:} Previous learning\par\textit{Available from 2026-09-03; suggested revisits:} 2026-09-10, 2026-09-24, 2026-11-03.\fi
\Needspace{5\baselineskip}\subsection*{R021 -- Units}\label{R021}
Convert 1 min 40 s to seconds; 2400 frames over this duration give how many per second?
\ifprof\textbf{Solution.} $100$ s; $24$ frames/s.\par\textit{Prerequisites:} Previous learning\par\textit{Available from 2026-09-03; suggested revisits:} 2026-09-10, 2026-09-24, 2026-11-03.\fi
\Needspace{5\baselineskip}\subsection*{R022 -- Simple probability}\label{R022}
With a fair die, probability of a multiple of 3?
\ifprof\textbf{Solution.} $\dfrac{2}{6}=\dfrac{1}{3}$.\par\textit{Prerequisites:} Previous learning\par\textit{Available from 2026-09-03; suggested revisits:} 2026-09-10, 2026-09-24, 2026-11-03.\fi
\Needspace{5\baselineskip}\subsection*{R023 -- Complement}\label{R023}
$P(A)=\dfrac{3}{8}$. Give $P(\overline A)$.
\ifprof\textbf{Solution.} $\dfrac{5}{8}$.\par\textit{Prerequisites:} Previous learning\par\textit{Available from 2026-09-03; suggested revisits:} 2026-09-10, 2026-09-24, 2026-11-03.\fi
\Needspace{5\baselineskip}\subsection*{R024 -- Interval}\label{R024}
Write $x\geq-2$ and $x<3$ as an interval.
\ifprof\textbf{Solution.} $[-2;3[$.\par\textit{Prerequisites:} Previous learning\par\textit{Available from 2026-09-03; suggested revisits:} 2026-09-10, 2026-09-24, 2026-11-03.\fi
\Needspace{5\baselineskip}\subsection*{R025 -- Logic}\label{R025}
Negate ``all the students passed.
\ifprof\textbf{Solution.} ``At least one student did not pass.\par\textit{Prerequisites:} Previous learning\par\textit{Available from 2026-09-03; suggested revisits:} 2026-09-10, 2026-09-24, 2026-11-03.\fi
\Needspace{5\baselineskip}\subsection*{R026 -- Histogram}\label{R026}
Two classes of widths 2 and 4 have the same count 20. In a histogram with areas proportional to counts, what are the relative heights?
\ifprof\textbf{Solution.} Heights 10 and 5: the area, not the height alone, encodes the count.\par\textit{Prerequisites:} Previous learning\par\textit{Available from 2026-09-03; suggested revisits:} 2026-09-10, 2026-09-24, 2026-11-03.\fi
\Needspace{5\baselineskip}\subsection*{R027 -- Pie chart}\label{R027}
A sector measures $72^\circ$. Proportion?
\ifprof\textbf{Solution.} $\dfrac{72}{360}=20\,\%$.\par\textit{Prerequisites:} Previous learning\par\textit{Available from 2026-09-03; suggested revisits:} 2026-09-10, 2026-09-24, 2026-11-03.\fi
\Needspace{5\baselineskip}\subsection*{R028 -- Scale}\label{R028}
An axis starts at 90. Bars stop at 100 and 110. Is the second value double the first?
\ifprof\textbf{Solution.} No: $\dfrac{110}{100}=1.1$. The heights drawn above 90 make this comparison misleading.\par\textit{Prerequisites:} Previous learning\par\textit{Available from 2026-09-03; suggested revisits:} 2026-09-10, 2026-09-24, 2026-11-03.\fi
\Needspace{5\baselineskip}\subsection*{R029 -- Percentage of a percentage}\label{R029}
40 \% of the students take A; among them, 25 \% take B. What proportion takes A and B?
\ifprof\textbf{Solution.} $0.4\times0.25=10\,\%$.\par\textit{Prerequisites:} Previous learning\par\textit{Available from 2026-09-03; suggested revisits:} 2026-09-10, 2026-09-24, 2026-11-03.\fi
\Needspace{5\baselineskip}\subsection*{R030 -- Proportion and count}\label{R030}
A category represents 0.125 of 80 people. Count?
\ifprof\textbf{Solution.} $10$ people.\par\textit{Prerequisites:} Previous learning\par\textit{Available from 2026-09-03; suggested revisits:} 2026-09-10, 2026-09-24, 2026-11-03.\fi
\Needspace{5\baselineskip}\subsection*{R031 -- Defining and representing a sequence}\label{R031}
Calculate the first four terms of $w_0=2$, $w_{n+1}=w_n+2n+1$; suggest an explicit formula.
\ifprof\textbf{Solution.} $2,3,6,11$. The formula $w_n=n^2+2$ works: it gives the first term and $(n+1)^2+2-(n^2+2)=2n+1$.\par\textit{Prerequisites:} Algebraic manipulation, integer powers, percentages, affine function, Python loop.\par\textit{Available from 2026-09-09; suggested revisits:} 2026-09-17, 2026-10-01, 2026-11-03.\fi
\Needspace{5\baselineskip}\subsection*{R032 -- Arithmetic sequences}\label{R032}
We know $u_3=17$ and $u_8=37$ for an arithmetic sequence. Determine $r$ and $u_0$.
\ifprof\textbf{Solution.} $37-17=5r$, so $r=4$ and then $u_0=17-3\times4=5$.\par\textit{Prerequisites:} Algebraic manipulation, integer powers, percentages, affine function, Python loop.\par\textit{Available from 2026-09-11; suggested revisits:} 2026-09-22, 2026-10-06, 2026-11-03.\fi
\Needspace{5\baselineskip}\subsection*{R033 -- Geometric sequences}\label{R033}
Study the variation of $u_n=-2\times3^n$, then of $v_n=2(-1)^n$.
\ifprof\textbf{Solution.} $u_{n+1}-u_n=-4\times3^n<0$: decreasing. $v_n$ alternates between 2 and $-2$: neither increasing nor decreasing.\par\textit{Prerequisites:} Algebraic manipulation, integer powers, percentages, affine function, Python loop.\par\textit{Available from 2026-09-16; suggested revisits:} 2026-09-24, 2026-10-08, 2026-11-05.\fi
\Needspace{5\baselineskip}\subsection*{R034 -- Sums, indices and proofs}\label{R034}
Calculate $\sum_{k=2}^{5}2^k$ without the summation symbol in the answer.
\ifprof\textbf{Solution.} $4+8+16+32=60$, or $\dfrac{4(1-2^4)}{1-2}=60$.\par\textit{Prerequisites:} Algebraic manipulation, integer powers, percentages, affine function, Python loop.\par\textit{Available from 2026-09-18; suggested revisits:} 2026-09-29, 2026-10-13, 2026-11-10.\fi
\Needspace{5\baselineskip}\subsection*{R035 -- Conjecturing a limit and programming a threshold}\label{R035}
For $u_n=48(\dfrac{1}{2})^n$, determine the first indices where $u_n\leq6$ and where $u_n<6$.
\ifprof\textbf{Solution.} $u_3=6$: index 3 for $\leq6$, index 4 for $<6$.\par\textit{Prerequisites:} Algebraic manipulation, integer powers, percentages, affine function, Python loop.\par\textit{Available from 2026-09-18; suggested revisits:} 2026-09-29, 2026-10-13, 2026-11-10.\fi
\Needspace{5\baselineskip}\subsection*{R036 -- Three forms, three readings}\label{R036}
Give three forms of $g(x)=x^2+2x-3$. Read off its roots and its minimum.
\ifprof\textbf{Solution.} $g(x)=(x+1)^2-4=(x-1)(x+3)$. Roots 1 and $-3$, minimum $-4$ at $-1$.\par\textit{Prerequisites:} Special products, zero-product property, sign of a product, square roots.\par\textit{Available from 2026-09-23; suggested revisits:} 2026-10-01, 2026-10-15, 2026-11-12.\fi
\Needspace{5\baselineskip}\subsection*{R037 -- Discriminant and solving}\label{R037}
Solve $3x^2-6x+3=0$ and $x^2+2x+5=0$.
\ifprof\textbf{Solution.} The first is $3(x-1)^2=0$, so 1. For the second, $\Delta=-16<0$: no real solution.\par\textit{Prerequisites:} Special products, zero-product property, sign of a product, square roots.\par\textit{Available from 2026-09-25; suggested revisits:} 2026-10-06, 2026-11-03, 2026-11-17.\fi
\Needspace{5\baselineskip}\subsection*{R038 -- Factorising without automating the discriminant}\label{R038}
Determine the polynomial f of degree 2 such that $f(1)=f(4)=0$ and $f(0)=-8$.
\ifprof\textbf{Solution.} $f(x)=a(x-1)(x-4)$; $4a=-8$, so $f(x)=-2(x-1)(x-4)$.\par\textit{Prerequisites:} Special products, zero-product property, sign of a product, square roots.\par\textit{Available from 2026-09-30; suggested revisits:} 2026-10-08, 2026-11-03, 2026-11-19.\fi
\Needspace{5\baselineskip}\subsection*{R039 -- Sign, inequalities and optimisation}\label{R039}
Solve $x^2-4x-5<0$, then maximise $x(8-x)$ on $[0;3]$.
\ifprof\textbf{Solution.} $(x-5)(x+1)<0$ gives $]-1;5[$. The vertex of $x(8-x)$ is at 4, outside the domain; on $[0;3]$ the function is increasing and its maximum is 15 at 3.\par\textit{Prerequisites:} Special products, zero-product property, sign of a product, square roots.\par\textit{Available from 2026-10-02; suggested revisits:} 2026-10-13, 2026-11-03, 2026-11-24.\fi
\Needspace{5\baselineskip}\subsection*{R040 -- Projection and angle}\label{R040}
Two vectors have norms 2 and 5 and form $90^\circ$, then $180^\circ$. Calculate the products.
\ifprof\textbf{Solution.} $0$ then $-10$. The cosine is 0 and $-1$ respectively.\par\textit{Prerequisites:} Vectors, Chasles's relation, Pythagoras, cosine in a right triangle, orthonormal frame.\par\textit{Available from 2026-10-07; suggested revisits:} 2026-10-15, 2026-11-03, 2026-11-26.\fi
\Needspace{5\baselineskip}\subsection*{R041 -- Coordinates, norm and orthogonality}\label{R041}
For which real number t are the vectors $(t;2)$ and $(3;-6)$ orthogonal?
\ifprof\textbf{Solution.} $3t-12=0$, so $t=4$.\par\textit{Prerequisites:} Vectors, Chasles's relation, Pythagoras, cosine in a right triangle, orthonormal frame.\par\textit{Available from 2026-10-09; suggested revisits:} 2026-11-03, 2026-11-03, 2026-12-01.\fi
\Needspace{5\baselineskip}\subsection*{R042 -- Vector calculus and identities}\label{R042}
Given $\norm{\vec u}=2$, $\norm{\vec v}=3$, $\vec u\cdot\vec v=1$. Calculate $\norm{\vec u-\vec v}$.
\ifprof\textbf{Solution.} The square is $4-2+9=11$; the norm is $\sqrt{11}$.\par\textit{Prerequisites:} Vectors, Chasles's relation, Pythagoras, cosine in a right triangle, orthonormal frame.\par\textit{Available from 2026-10-14; suggested revisits:} 2026-11-03, 2026-11-05, 2026-12-03.\fi
\Needspace{5\baselineskip}\subsection*{R043 -- Al-Kashi and choosing a method}\label{R043}
A triangle has $AB=AC=4$ and $BC=4\sqrt2$. Determine the angle A.
\ifprof\textbf{Solution.} $32=16+16-32\cos A$, so $\cos A=0$ and $A=90^\circ$.\par\textit{Prerequisites:} Vectors, Chasles's relation, Pythagoras, cosine in a right triangle, orthonormal frame.\par\textit{Available from 2026-10-16; suggested revisits:} 2026-11-03, 2026-11-10, 2026-12-08.\fi
\Needspace{5\baselineskip}\subsection*{R044 -- The circle with diameter AB}\label{R044}
For $A(0;0)$, $B(6;0)$, determine the points M such that $\vect{MA}\cdot\vect{MB}=7$.
\ifprof\textbf{Solution.} $MI^2-9=7$: circle with centre $(3;0)$ and radius 4.\par\textit{Prerequisites:} Vectors, Chasles's relation, Pythagoras, cosine in a right triangle, orthonormal frame.\par\textit{Available from 2026-10-16; suggested revisits:} 2026-11-03, 2026-11-10, 2026-12-08.\fi
\Needspace{5\baselineskip}\subsection*{R045 -- Rate of change and secants}\label{R045}
Calculate the rate of $f(x)=x^2+3x$ between 2 and $2+h$.
\ifprof\textbf{Solution.} $f(2+h)-f(2)=7h+h^2$; rate $7+h$ for $h\ne0$.\par\textit{Prerequisites:} Images, slope of a line, algebraic manipulation, square root and absolute value.\par\textit{Available from 2026-11-04; suggested revisits:} 2026-11-12, 2026-11-26, 2027-01-05.\fi
\Needspace{5\baselineskip}\subsection*{R046 -- Derivative at a point}\label{R046}
For $f(x)=3x^2-2x$, calculate $f\prime(1)$ from the rate.
\ifprof\textbf{Solution.} The rate is $4+3h$, so the derivative at 1 is 4.\par\textit{Prerequisites:} Images, slope of a line, algebraic manipulation, square root and absolute value.\par\textit{Available from 2026-11-06; suggested revisits:} 2026-11-17, 2026-12-01, 2027-01-05.\fi
\Needspace{5\baselineskip}\subsection*{R047 -- Tangent and affine approximation}\label{R047}
Given $f(3)=5$ and $f\prime(3)=-2$. Write the tangent and approximate $f(3.04)$.
\ifprof\textbf{Solution.} $y=5-2(x-3)=-2x+11$; approximation $5-2(0.04)=4.92$.\par\textit{Prerequisites:} Images, slope of a line, algebraic manipulation, square root and absolute value.\par\textit{Available from 2026-11-11; suggested revisits:} 2026-11-19, 2026-12-03, 2027-01-05.\fi
\Needspace{5\baselineskip}\subsection*{R048 -- When differentiability fails}\label{R048}
Justify that $g(x)=|x-2|$ is not differentiable at 2. Also study the differentiability of $x\mapsto\sqrt{x}$ at 0 through its rate.
\ifprof\textbf{Solution.} The rate at 2 is $\dfrac{|h|}{h}$, with two opposite limiting values. For the square root at 0, the rate $\dfrac{1}{\sqrt h}$ for $h>0$ exceeds every bound as h gets closer to 0: no real derivative.\par\textit{Prerequisites:} Images, slope of a line, algebraic manipulation, square root and absolute value.\par\textit{Available from 2026-11-13; suggested revisits:} 2026-11-24, 2026-12-08, 2027-01-05.\fi
\Needspace{5\baselineskip}\subsection*{R049 -- Conditioning: changing the denominator}\label{R049}
Out of 60 objects, 20 are red; among the red ones, 5 are damaged. Calculate $P(R\cap D)$ and $P_R(D)$.
\ifprof\textbf{Solution.} $\dfrac{5}{60}=\dfrac{1}{12}$ and $\dfrac{5}{20}=\dfrac{1}{4}$.\par\textit{Prerequisites:} Fractions, proportions, probabilities on a finite sample space, intersection and union.\par\textit{Available from 2026-11-18; suggested revisits:} 2026-11-26, 2026-12-10, 2027-01-07.\fi
\Needspace{5\baselineskip}\subsection*{R050 -- Trees, partitions and total probability}\label{R050}
$P(A)=\dfrac{1}{3}$, $P_A(B)=\dfrac{3}{4}$, $P_{\overline A}(B)=\dfrac{1}{4}$. Calculate $P(B)$.
\ifprof\textbf{Solution.} $\dfrac{1}{4}+\dfrac{1}{6}=\dfrac{5}{12}$.\par\textit{Prerequisites:} Fractions, proportions, probabilities on a finite sample space, intersection and union.\par\textit{Available from 2026-11-20; suggested revisits:} 2026-12-01, 2026-12-15, 2027-01-12.\fi
\Needspace{5\baselineskip}\subsection*{R051 -- Independence and disjointness}\label{R051}
Given $P(A)=0.4$, $P(B)=0.5$, $P(A\cap B)=0.18$. Are they independent?
\ifprof\textbf{Solution.} No: $0.4\times0.5=0.20\ne0.18$.\par\textit{Prerequisites:} Fractions, proportions, probabilities on a finite sample space, intersection and union.\par\textit{Available from 2026-11-25; suggested revisits:} 2026-12-03, 2026-12-17, 2027-01-19.\fi
\Needspace{5\baselineskip}\subsection*{R052 -- Successive trials and Bernoulli trials}\label{R052}
Four independent trials, each with probability $p=\dfrac{1}{3}$. Calculate the probability of four failures, then of at least one success.
\ifprof\textbf{Solution.} $(\dfrac{2}{3})^4=\dfrac{16}{81}$ then $1-\dfrac{16}{81}=\dfrac{65}{81}$.\par\textit{Prerequisites:} Fractions, proportions, probabilities on a finite sample space, intersection and union.\par\textit{Available from 2026-11-27; suggested revisits:} 2026-12-08, 2027-01-05, 2027-01-19.\fi
\Needspace{5\baselineskip}\subsection*{R053 -- Simulating and discussing the model}\label{R053}
A simulation of 2000 points gives 675 points below $y=x^2$. What area does it estimate?
\ifprof\textbf{Solution.} $\dfrac{675}{2000}=0.3375$; this is an estimate, not the exact value.\par\textit{Prerequisites:} Fractions, proportions, probabilities on a finite sample space, intersection and union.\par\textit{Available from 2026-11-27; suggested revisits:} 2026-12-08, 2027-01-05, 2027-01-19.\fi
\Needspace{5\baselineskip}\subsection*{R054 -- Derivative function and standard functions}\label{R054}
Differentiate $x^5$, $x^{-2}$ and $\sqrt x$, stating the domains.
\ifprof\textbf{Solution.} $5x^4$ on $\R$; $-\dfrac{2}{x^3}$ for $x\ne0$; $\dfrac{1}{2\sqrt x}$ for $x>0$.\par\textit{Prerequisites:} Derivative at a point, algebraic calculation, sign of a quadratic, intervals.\par\textit{Available from 2026-12-02; suggested revisits:} 2026-12-10, 2027-01-05, 2027-01-21.\fi
\Needspace{5\baselineskip}\subsection*{R055 -- Sum, product, reciprocal and quotient}\label{R055}
Differentiate $(x^2+1)(3x-2)$, then $\dfrac{1}{x^2+1}$.
\ifprof\textbf{Solution.} $2x(3x-2)+3(x^2+1)=9x^2-4x+3$; $-\dfrac{2x}{(x^2+1)^2}$, on $\R$.\par\textit{Prerequisites:} Derivative at a point, algebraic calculation, sign of a quadratic, intervals.\par\textit{Available from 2026-12-04; suggested revisits:} 2026-12-15, 2027-01-05, 2027-01-26.\fi
\Needspace{5\baselineskip}\subsection*{R056 -- Sign of the derivative and extrema}\label{R056}
Study $f(x)=x^2-4x+1$ on $[0;3]$ and give its extrema on this interval.
\ifprof\textbf{Solution.} $f\prime=2x-4$, so decreasing up to 2 then increasing. Minimum $-3$ at 2; values at the endpoints 1 and $-2$, so maximum 1 at 0.\par\textit{Prerequisites:} Derivative at a point, algebraic calculation, sign of a quadratic, intervals.\par\textit{Available from 2026-12-09; suggested revisits:} 2026-12-17, 2027-01-05, 2027-01-28.\fi
\Needspace{5\baselineskip}\subsection*{R057 -- Parity, inequalities and comparison}\label{R057}
Study the parity of $x^4+x^2$ and $x^3-x$; prove $x^2+1\geq2x$.
\ifprof\textbf{Solution.} First even, second odd; difference $x^2+1-2x=(x-1)^2\geq0$, equality at 1.\par\textit{Prerequisites:} Derivative at a point, algebraic calculation, sign of a quadratic, intervals.\par\textit{Available from 2026-12-11; suggested revisits:} 2027-01-05, 2027-01-05, 2027-02-02.\fi
\Needspace{5\baselineskip}\subsection*{R058 -- Optimisation and Newton's method}\label{R058}
A rectangle of perimeter 12 m has width x. Determine the maximum area by differentiation.
\ifprof\textbf{Solution.} $A(x)=x(6-x)$ on $[0;6]$; $A\prime=6-2x$. Maximum 9 m$^2$ at 3.\par\textit{Prerequisites:} Derivative at a point, algebraic calculation, sign of a quadratic, intervals.\par\textit{Available from 2026-12-11; suggested revisits:} 2027-01-05, 2027-01-05, 2027-02-02.\fi
\Needspace{5\baselineskip}\subsection*{R059 -- Normal vector and Cartesian equation}\label{R059}
Equation of the line through $(2;-1)$ with normal $(3;4)$.
\ifprof\textbf{Solution.} $3(x-2)+4(y+1)=0$, that is $3x+4y-2=0$.\par\textit{Prerequisites:} Vector coordinates, scalar product, linear systems, completing the square. The frame is always orthonormal.\par\textit{Available from 2027-01-06; suggested revisits:} 2027-01-19, 2027-01-28, 2027-03-02.\fi
\Needspace{5\baselineskip}\subsection*{R060 -- Orthogonal projection and distance}\label{R060}
Project $P(4;1)$ onto $d:x=2$.
\ifprof\textbf{Solution.} $H(2;1)$ and distance 2; the normal is horizontal.\par\textit{Prerequisites:} Vector coordinates, scalar product, linear systems, completing the square. The frame is always orthonormal.\par\textit{Available from 2027-01-08; suggested revisits:} 2027-01-19, 2027-02-02, 2027-03-02.\fi
\Needspace{5\baselineskip}\subsection*{R061 -- Equation of a circle}\label{R061}
Recognise $x^2+y^2+2x-8y+17=0$.
\ifprof\textbf{Solution.} $(x+1)^2+(y-4)^2=0$: the single point $(-1;4)$, not a circle of strictly positive radius.\par\textit{Prerequisites:} Vector coordinates, scalar product, linear systems, completing the square. The frame is always orthonormal.\par\textit{Available from 2027-01-13; suggested revisits:} 2027-01-21, 2027-02-04, 2027-03-04.\fi
\Needspace{5\baselineskip}\subsection*{R062 -- Solving a configuration}\label{R062}
Intersect $(x-1)^2+(y+2)^2=9$ with $y=1$.
\ifprof\textbf{Solution.} $(x-1)^2=0$, so the only point is $(1;1)$; the line is tangent to the circle.\par\textit{Prerequisites:} Vector coordinates, scalar product, linear systems, completing the square. The frame is always orthonormal.\par\textit{Available from 2027-01-15; suggested revisits:} 2027-01-26, 2027-02-09, 2027-03-09.\fi
\Needspace{5\baselineskip}\subsection*{R063 -- Variable, event and distribution}\label{R063}
Two fair coins are tossed. X counts the heads. Write its distribution and $P(X\leq1)$.
\ifprof\textbf{Solution.} Values 0, 1, 2; probabilities $\dfrac{1}{4}$, $\dfrac{1}{2}$, $\dfrac{1}{4}$. $P(X\leq1)=\dfrac{3}{4}$.\par\textit{Prerequisites:} Probability, weighted mean, squares, trees and lists.\par\textit{Available from 2027-01-20; suggested revisits:} 2027-01-28, 2027-02-11, 2027-03-11.\fi
\Needspace{5\baselineskip}\subsection*{R064 -- Expectation and linearity}\label{R064}
A gross gain is 10 with probability $\dfrac{1}{5}$ and 0 otherwise. What stake makes the game fair?
\ifprof\textbf{Solution.} $E(G)=\dfrac{10}{5}=2$. For a stake m, $E(G-m)=2-m$; fair stake 2.\par\textit{Prerequisites:} Probability, weighted mean, squares, trees and lists.\par\textit{Available from 2027-01-22; suggested revisits:} 2027-02-02, 2027-03-02, 2027-03-16.\fi
\Needspace{5\baselineskip}\subsection*{R065 -- Variance, standard deviation and König--Huygens}\label{R065}
For X with values $-1,1$ equally likely, calculate the expectation, variance and standard deviation.
\ifprof\textbf{Solution.} $E(X)=0$, $E(X^2)=1$, variance 1 and standard deviation 1.\par\textit{Prerequisites:} Probability, weighted mean, squares, trees and lists.\par\textit{Available from 2027-01-27; suggested revisits:} 2027-02-04, 2027-03-02, 2027-03-18.\fi
\Needspace{5\baselineskip}\subsection*{R066 -- Transforming and comparing gains}\label{R066}
$E(X)=3$, $V(X)=4$. Give the characteristics of $Y=-2X+5$.
\ifprof\textbf{Solution.} $E(Y)=-1$, $V(Y)=16$, $\sigma(Y)=4$.\par\textit{Prerequisites:} Probability, weighted mean, squares, trees and lists.\par\textit{Available from 2027-01-29; suggested revisits:} 2027-02-09, 2027-03-02, 2027-03-23.\fi
\Needspace{5\baselineskip}\subsection*{R067 -- Samples and observed mean}\label{R067}
For $\mu=5$, $\sigma=3$ and n=36, give the experimental interval asked for.
\ifprof\textbf{Solution.} $\dfrac{2\sigma}{\sqrt n}=1$; interval $[4;6]$.\par\textit{Prerequisites:} Probability, weighted mean, squares, trees and lists.\par\textit{Available from 2027-01-29; suggested revisits:} 2027-02-09, 2027-03-02, 2027-03-23.\fi
\Needspace{5\baselineskip}\subsection*{R068 -- Defining a new function}\label{R068}
What is the equation of the tangent to $e^x$ at 1?
\ifprof\textbf{Solution.} $y=e+e(x-1)=ex$.\par\textit{Prerequisites:} Powers, geometric sequences, differentiation and sign.\par\textit{Available from 2027-02-03; suggested revisits:} 2027-02-11, 2027-03-02, 2027-03-25.\fi
\Needspace{5\baselineskip}\subsection*{R069 -- Algebraic properties and comparisons}\label{R069}
Simplify $\dfrac{e^{3x}}{e^{x-1}}$; solve $e^{x+2}>e^{2x}$.
\ifprof\textbf{Solution.} $e^{2x+1}$; $x+2>2x$, so $x<2$.\par\textit{Prerequisites:} Powers, geometric sequences, differentiation and sign.\par\textit{Available from 2027-02-05; suggested revisits:} 2027-03-02, 2027-03-02, 2027-03-30.\fi
\Needspace{5\baselineskip}\subsection*{R070 -- Differentiating and studying}\label{R070}
Differentiate $g(t)=3e^{-2t}$ and study its direction of variation.
\ifprof\textbf{Solution.} $g\prime(t)=-6e^{-2t}<0$ on $\R$: strictly decreasing.\par\textit{Prerequisites:} Powers, geometric sequences, differentiation and sign.\par\textit{Available from 2027-02-10; suggested revisits:} 2027-03-02, 2027-03-04, 2027-04-06.\fi
\Needspace{5\baselineskip}\subsection*{R071 -- From discrete to continuous}\label{R071}
For $Q(t)=5e^{0.3t}$, express $Q(n)$ and $\dfrac{Q(t+2)}{Q(t)}$.
\ifprof\textbf{Solution.} $u_n=5(e^{0.3})^n$; the ratio over two units is $e^{0.6}$.\par\textit{Prerequisites:} Powers, geometric sequences, differentiation and sign.\par\textit{Available from 2027-02-12; suggested revisits:} 2027-03-02, 2027-03-09, 2027-04-06.\fi
\Needspace{5\baselineskip}\subsection*{R072 -- Building numerically and proving an inequality}\label{R072}
Show that $e^{2x}\geq1+2x$ for every real number x.
\ifprof\textbf{Solution.} Apply the inequality proved with the argument $2x$; equality only at x=0.\par\textit{Prerequisites:} Powers, geometric sequences, differentiation and sign.\par\textit{Available from 2027-02-12; suggested revisits:} 2027-03-02, 2027-03-09, 2027-04-06.\fi
\Needspace{5\baselineskip}\subsection*{R073 -- Measuring an angle in radians}\label{R073}
Convert $45^\circ$, $120^\circ$; give the arc for radius 2 m and angle $\dfrac{3\pi}{4}$.
\ifprof\textbf{Solution.} $\dfrac{\pi}{4}$, $\dfrac{2\pi}{3}$; length $\dfrac{3\pi}{2}$ m.\par\textit{Prerequisites:} Circle, Pythagoras, cosine and sine in a right triangle, symmetries, fractions.\par\textit{Available from 2027-03-03; suggested revisits:} 2027-03-11, 2027-03-25, 2027-04-27.\fi
\Needspace{5\baselineskip}\subsection*{R074 -- Wrapping and the image of a real number}\label{R074}
Place the images of $\dfrac{7\pi}{2}$ and $-\dfrac{5\pi}{3}$ by reducing them to within one turn.
\ifprof\textbf{Solution.} $\dfrac{7\pi}{2}$ has the image of $\dfrac{3\pi}{2}$; $-\dfrac{5\pi}{3}$ that of $\dfrac{\pi}{3}$.\par\textit{Prerequisites:} Circle, Pythagoras, cosine and sine in a right triangle, symmetries, fractions.\par\textit{Available from 2027-03-05; suggested revisits:} 2027-03-16, 2027-03-30, 2027-04-27.\fi
\Needspace{5\baselineskip}\subsection*{R075 -- Sine, cosine and standard values}\label{R075}
If $\cos x=\dfrac{3}{5}$ and the image of x is in the fourth quadrant, give $\sin x$.
\ifprof\textbf{Solution.} $\sin ^2x=\dfrac{16}{25}$; the sine is negative, so $-\dfrac{4}{5}$.\par\textit{Prerequisites:} Circle, Pythagoras, cosine and sine in a right triangle, symmetries, fractions.\par\textit{Available from 2027-03-10; suggested revisits:} 2027-03-18, 2027-04-06, 2027-04-29.\fi
\Needspace{5\baselineskip}\subsection*{R076 -- Related angles and symmetries}\label{R076}
Calculate the two coordinates of the images of $\dfrac{4\pi}{3}$ and $\dfrac{7\pi}{4}$.
\ifprof\textbf{Solution.} $(-\dfrac{1}{2};-\dfrac{\sqrt 3}{2})$ and $(\dfrac{\sqrt 2}{2};-\dfrac{\sqrt 2}{2})$.\par\textit{Prerequisites:} Circle, Pythagoras, cosine and sine in a right triangle, symmetries, fractions.\par\textit{Available from 2027-03-12; suggested revisits:} 2027-03-23, 2027-04-06, 2027-05-04.\fi
\Needspace{5\baselineskip}\subsection*{R077 -- Approximating pi and putting it to use}\label{R077}
For n=6 and c=1, calculate the new chord after doubling.
\ifprof\textbf{Solution.} $h=\dfrac{\sqrt 3}{2}$; new chord $\dfrac{1}{\sqrt {2+\sqrt3}}=\sqrt {2-\sqrt3}$.\par\textit{Prerequisites:} Circle, Pythagoras, cosine and sine in a right triangle, symmetries, fractions.\par\textit{Available from 2027-03-12; suggested revisits:} 2027-03-23, 2027-04-06, 2027-05-04.\fi
\Needspace{20\baselineskip}
\Needspace{5\baselineskip}\subsection*{R078 -- Reading a graph}\label{R078}
On the figure R078, read $f(1)$, solve $f(x)=0$ and $f(x)>0$.
\begin{center}\begin{tikzpicture}\begin{axis}[width=.8\linewidth,height=5cm,axis lines=middle,xmin=-2.5,xmax=4.5,ymin=-2,ymax=3,xtick={-2,-1,0,1,2,3,4},ytick={-1,0,1,2},grid=both]\addplot[very thick,vert,mark=*] coordinates {(-2,-1)(1,2)(4,-1)};\end{axis}\end{tikzpicture}\end{center}\ifprof\textbf{Solution.} $f(1)=2$; zeros $-1$ and 3; positive sign on $]-1;3[$.\par\textit{Prerequisites:} Plotting points\par\textit{Available from 2026-09-03; suggested revisits:} 2026-09-10, 2026-09-24, 2026-11-03.\fi
